
\medskip
\noindent\SGARefTarget{sga3:xi:remark:1:6}\textbf{Remark 1.6.} ---
Of course, given a morphism \(f:X\to Y\) of preschemes, one says that
\(f\) is smooth (respectively, unramified, respectively, étale) if
\(f\) makes \(X\) a smooth (respectively, unramified, respectively,
étale) \(Y\)-prescheme.  When \(X\) and \(Y\) are \(S\)-preschemes and
\(f\) is an \(S\)-morphism of finite presentation, these properties of
\(f\) can be expressed immediately in terms of the morphism of
functors
\[
  F,G:(\mathbf{Sch})^\circ/S\longrightarrow(\mathbf{Ens})
\]
defined by \(f\).  Namely, \(f\) is smooth (respectively, unramified,
respectively, étale) if, for every \(S\)-prescheme \(S'\) that is
affine in the absolute sense and every subscheme \(S'_0\subset S'\)
defined by a nilpotent ideal \(\mathscr J\), the map
\[
  F(S')\longrightarrow
  F(S'_0)\times_{G(S'_0)}G(S')
\]
deduced from the commutative square

\SGARefTarget{sga3:xi:diagram:1:6:functor-square}
\[
\begin{tikzcd}[column sep=huge,row sep=large]
  F(S') \arrow[r] \arrow[d] & G(S') \arrow[d] \\
  F(S'_0) \arrow[r] & G(S'_0)
\end{tikzcd}
\]
is surjective (respectively, injective, respectively, bijective).
Equivalently, for every commutative square of morphisms of functors
over \(S\), where \(S'\) and \(S'_0\) are as above and
\(i:S'_0\to S'\) is the canonical immersion,

\SGARefTarget{sga3:xi:diagram:1:6:lifting-square}
\[
\begin{tikzcd}[column sep=huge,row sep=large]
  \text{\rm(Q)} & S'_0 \arrow[r,"i"] \arrow[d,"u_0"] &
    S' \arrow[d,"v"'] \\
  & F \arrow[r,"f"] & G ,
\end{tikzcd}
\]
there exists at least one (respectively, at most one, respectively,
exactly one) morphism

\SGARefTarget{sga3:xi:diagram:1:6:lift}
\[
\begin{tikzcd}[column sep=huge,row sep=large]
  & S' \arrow[dl,"u"'] \\
  F &
\end{tikzcd}
\]
that makes the two corresponding triangles commute:
\[
  u_0=u\circ i
  \qquad\text{and}\qquad
  v=f\circ u.
\]

This property for a homomorphism of functors \(f:F\to G\) over \(S\)
retains its meaning even when the functors under consideration are
not representable.  If it holds, the homomorphism \(f:F\to G\) is
said to be \emph{formally smooth} (respectively,
\emph{formally unramified}, respectively, \emph{formally étale}).
Notice that this property depends only on the homomorphism of the
underlying functors
\[
  (\mathbf{Sch})^\circ\longrightarrow(\mathbf{Ens})
\]
defined by \(F\) and \(G\) (cf.\ \SGARefLink{sga3:i:proposition:1:4:1}{Exposé I, 1.4.1}), and not on the
structure morphisms \(F\to S\) and \(G\to S\).  An equivalent
formulation of the preceding definition, somewhat easier to use in
applications, is the following: a morphism of functors \(F\to G\)
over \(S\) is formally smooth (respectively, formally unramified,
respectively, formally étale) if, for every \(S'\) over \(S\) and
every morphism \(S'\to G\), the functor
\[
  F'=F\times_G S'
\]
over \(S'\) is formally smooth (respectively, formally unramified,
respectively, formally étale).
